% 第五章至第七章
% 建议拆分路径：chapters/6_chapter6.tex
% 写作说明：本内容基于论文目录大纲第6章扩展。
% 风格说明：学术表达约 70%，自然化说明约 30%；避免重复前文已解释的技术概念。
% 图表说明：本文中的 figure 为可编译示意图，你也可以后续替换为正式图片。

\chapter{系统测试与结果分析}
\label{chap:system-test-analysis}

系统测试的目的是验证在线视频分享平台的核心功能是否能够按照需求正常运行，并检查微服务调用、视频上传转码、前端交互、后台审核和 AI 审核等关键流程是否具备完整性。本章主要从测试环境、测试方案、核心功能测试、AI 审核测试、异常场景测试以及测试结果分析几个方面进行说明。

由于本系统涉及多个服务和中间件，测试不能只关注单个接口是否可用，还需要关注跨服务业务链路是否能够正确流转。例如，视频上传成功并不代表投稿流程完成，还需要继续验证转码是否成功、审核任务是否创建、AI 结果是否回写、管理员是否能审核以及前台是否能播放。软件质量模型中将功能适合性、性能效率、可靠性、安全性和可维护性都视为软件质量的重要方面，因此测试也应围绕这些方面展开\cite{iso25010}。

\section{测试环境与测试方案}
\label{sec:test-env-plan}

\subsection{测试环境配置}
\label{subsec:test-env-config}

系统测试环境主要包括前端运行环境、Java 后端运行环境、Python AI 服务环境以及数据库和中间件环境。测试前需要启动 MySQL、Redis、Elasticsearch、RabbitMQ、Nacos 和 Seata 等基础组件，再启动后端网关、主站服务、后台管理服务、资源服务、互动服务和 AI 审核服务，最后启动用户端和管理端页面。

表~\ref{tab:test-env-config} 展示了系统测试环境配置。该表可以看出，本系统测试涉及的组件较多，因此测试前需要确认各组件是否正常运行。

\begin{table}[htbp]
  \centering
  \caption{系统测试环境配置}
  \label{tab:test-env-config}
  \begin{tabular}{p{0.22\textwidth}p{0.30\textwidth}p{0.38\textwidth}}
    \hline
    \textbf{环境类别} & \textbf{配置项} & \textbf{说明} \\
    \hline
    操作系统 & Windows / Linux & 本地开发与演示环境均可运行 \\
    后端环境 & JDK 17、Maven & 用于运行 Spring Boot / Spring Cloud 微服务 \\
    AI 服务环境 & Python、CUDA 环境 & 用于运行 AI 审核服务和模型推理 \\
    数据库 & MySQL 8 & 保存用户、视频、投稿、互动和审核数据 \\
    缓存 & Redis & 保存登录态、验证码、上传状态和队列数据 \\
    搜索引擎 & Elasticsearch & 支持视频全文检索 \\
    消息队列 & RabbitMQ & 传递 AI 审核请求和审核结果 \\
    注册中心 & Nacos & 管理微服务注册发现和配置 \\
    视频工具 & FFmpeg & 完成视频转码、切片和音频提取 \\
    \hline
  \end{tabular}
\end{table}

由表~\ref{tab:test-env-config} 可知，系统测试依赖多个基础组件。如果某个组件未启动或配置错误，可能会导致部分功能无法正常运行。例如 RabbitMQ 异常会影响 AI 审核任务投递，FFmpeg 不可用会导致视频无法转码，Elasticsearch 异常会影响视频搜索。

\subsection{测试方案设计}
\label{subsec:test-plan-design}

本系统测试采用功能测试、流程测试和异常测试相结合的方式。功能测试用于验证单个模块是否满足需求，例如登录、视频列表、评论、弹幕和分类管理；流程测试用于验证完整业务链路是否顺畅，例如视频从上传到转码、AI 审核、人工审核再到前台播放的过程；异常测试用于验证系统在上传失败、转码失败、AI 审核失败、未登录访问和权限不足等情况下是否能够给出合理处理。

测试方案如表~\ref{tab:test-plan} 所示。该表把测试类型、测试内容和测试目标放在一起，有助于保证测试覆盖系统主要功能。

\begin{table}[htbp]
  \centering
  \caption{系统测试方案}
  \label{tab:test-plan}
  \begin{tabular}{p{0.24\textwidth}p{0.38\textwidth}p{0.28\textwidth}}
    \hline
    \textbf{测试类型} & \textbf{测试内容} & \textbf{测试目标} \\
    \hline
    用户端功能测试 & 首页、搜索、播放、评论、弹幕、点赞、收藏、投币 & 验证普通用户功能可正常使用 \\
    创作中心测试 & 封面上传、视频分片上传、投稿提交、稿件状态查看 & 验证创作者投稿流程完整 \\
    后台管理测试 & 管理员登录、分类管理、视频审核 & 验证后台管理功能可用 \\
    视频处理测试 & 分片保存、转码、HLS 播放、播放统计 & 验证资源服务处理能力 \\
    AI 审核测试 & 审核任务创建、消息发送、结果回写、后台展示 & 验证 AI 审核闭环 \\
    异常场景测试 & 未登录操作、转码失败、AI 失败、重复审核 & 验证系统异常处理能力 \\
    \hline
  \end{tabular}
\end{table}

由表~\ref{tab:test-plan} 可知，测试范围覆盖了用户端、创作中心、后台管理、视频资源处理和 AI 审核等核心模块。这样的测试方式能更接近真实使用场景，而不是只验证单个接口是否返回成功。

为了验证系统是否形成完整的视频业务闭环，测试不能只检查单个接口是否返回成功，还需要从视频投稿开始，一直验证到视频发布和前台播放。算法~\ref{alg:test-full-video-flow} 给出了视频投稿到发布的核心链路测试流程。

\begin{algorithm}[htbp]
  \caption{视频投稿到发布的核心链路测试流程}
  \label{alg:test-full-video-flow}
  \begin{algorithmic}[1]
    \Require 测试用户账号，测试视频文件，测试管理员账号
    \Ensure 核心链路测试结果
    \State 使用测试用户登录系统
    \State 调用预上传接口生成上传任务
    \State 按分片上传测试视频文件
    \State 提交视频投稿信息
    \State 检查投稿记录状态是否为“转码中”
    \State 等待资源服务消费转码队列
    \State 检查分 P 文件是否生成 HLS 播放资源
    \State 检查投稿状态是否变为“待审核”
    \State 检查 AI 审核任务是否创建并发送消息
    \State 等待 AI 审核结果回写
    \State 使用管理员账号登录后台
    \State 查看审核列表和 AI 审核结果
    \State 执行人工审核通过
    \State 检查正式视频表和搜索索引是否写入成功
    \State 在前台搜索并播放该视频
    \If{上述步骤均符合预期}
      \State 测试通过
    \Else
      \State 记录失败步骤和异常信息
    \EndIf
  \end{algorithmic}
\end{algorithm}

由算法~\ref{alg:test-full-video-flow} 可知，核心链路测试覆盖了上传、转码、AI 审核、人工审核、正式发布、搜索和播放等多个环节。该测试方式比单独测试某个接口更能反映系统整体业务流程是否正确。

\section{核心功能测试}
\label{sec:core-function-test}

\subsection{用户认证与首页测试}
\label{subsec:user-auth-home-test}

用户认证测试主要验证注册、登录、退出和自动登录等功能是否正常。测试时输入正确账号和密码，应能够登录成功并进入用户端；输入错误密码或验证码，应返回明确错误提示；退出登录后，用户不应继续执行需要身份的操作。

首页测试主要验证视频列表加载、分类筛选和关键词搜索是否正常。系统应只展示审核通过并正式发布的视频，未审核或审核失败的视频不应出现在首页列表中。搜索时，用户输入关键词后，系统应返回与标题、标签或简介相关的视频内容。

表~\ref{tab:user-auth-home-test} 展示了用户认证与首页测试用例。该表说明了测试输入、预期结果和测试结论。

\begin{table}[htbp]
  \centering
  \caption{用户认证与首页测试用例}
  \label{tab:user-auth-home-test}
  \begin{tabular}{p{0.22\textwidth}p{0.32\textwidth}p{0.34\textwidth}}
    \hline
    \textbf{测试项} & \textbf{测试操作} & \textbf{预期结果} \\
    \hline
    用户登录 & 输入正确账号、密码和验证码 & 登录成功，进入用户端页面 \\
    错误登录 & 输入错误密码或验证码 & 返回错误提示，不允许登录 \\
    退出登录 & 点击退出按钮 & 清除登录态，限制互动操作 \\
    首页列表 & 打开首页 & 展示审核通过的视频列表 \\
    分类筛选 & 选择视频分类 & 返回对应分类下的视频 \\
    关键词搜索 & 输入关键词搜索 & 返回相关视频结果 \\
    \hline
  \end{tabular}
\end{table}

测试结果表明，用户认证和首页功能能够满足基本需求。登录态能够影响互动和投稿等操作，首页列表也能按正式发布状态展示视频。这里最重要的是保证未发布视频不会被普通用户看到。

\subsection{视频播放与互动测试}
\label{subsec:video-play-interact-test}

视频播放测试主要验证视频详情加载、HLS 播放资源访问和分 P 切换是否正常。用户进入播放页后，系统应加载视频基本信息、作者信息、分 P 文件和播放资源。点击不同分 P 时，播放器应切换到对应视频资源。

互动测试主要包括评论、弹幕、点赞、收藏、投币和关注。登录用户执行互动操作后，系统应保存对应行为，并在页面上更新状态；未登录用户执行互动操作时，系统应提示登录。弹幕测试还需要关注播放时间点，发送后的弹幕应能在对应时间附近展示。

表~\ref{tab:play-interact-test} 展示了视频播放与互动测试用例。

\begin{table}[htbp]
  \centering
  \caption{视频播放与互动测试用例}
  \label{tab:play-interact-test}
  \begin{tabular}{p{0.22\textwidth}p{0.34\textwidth}p{0.32\textwidth}}
    \hline
    \textbf{测试项} & \textbf{测试操作} & \textbf{预期结果} \\
    \hline
    视频详情 & 打开视频播放页 & 展示标题、简介、作者和统计信息 \\
    视频播放 & 加载 HLS 播放资源 & 视频能够正常播放 \\
    分 P 切换 & 点击不同分 P & 播放资源切换正常 \\
    评论发送 & 登录用户发表评论 & 评论保存并展示在评论区 \\
    弹幕发送 & 登录用户发送弹幕 & 弹幕在对应播放时间展示 \\
    用户行为 & 点赞、收藏、投币、关注 & 操作状态和统计数据更新 \\
    未登录互动 & 游客尝试互动 & 提示登录，不保存行为 \\
    \hline
  \end{tabular}
\end{table}

测试结果表明，播放页能够完成视频播放和主要互动功能。对于未登录用户，系统能够进行限制；对于登录用户，系统能够记录互动行为并返回状态。播放测试中需要特别关注资源路径是否正确，因为 HLS 播放依赖 \texttt{.m3u8} 和媒体切片文件。

\subsection{创作中心与视频上传测试}
\label{subsec:creator-upload-test}

创作中心测试重点验证视频投稿流程是否完整。测试时，创作者先上传封面，再选择视频文件进行分片上传，最后填写标题、分类、标签和简介并提交投稿。系统应保存投稿主数据和分 P 文件数据，并将文件加入转码队列。

视频上传测试需要关注分片是否完整保存，上传状态是否正确记录，分片合并是否成功。如果上传过程中某个分片失败，系统应允许重新上传失败分片，而不是要求用户重新上传整个视频。

表~\ref{tab:creator-upload-test} 展示了创作中心测试用例。

\begin{table}[htbp]
  \centering
  \caption{创作中心与视频上传测试用例}
  \label{tab:creator-upload-test}
  \begin{tabular}{p{0.22\textwidth}p{0.34\textwidth}p{0.32\textwidth}}
    \hline
    \textbf{测试项} & \textbf{测试操作} & \textbf{预期结果} \\
    \hline
    封面上传 & 上传视频封面图片 & 图片保存成功，返回访问路径 \\
    预上传 & 选择视频并申请上传任务 & 返回 uploadId 和上传状态 \\
    分片上传 & 按分片顺序上传视频文件 & 分片保存成功，进度更新 \\
    投稿提交 & 填写标题、分类、标签和简介 & 投稿数据保存成功 \\
    状态查看 & 查看稿件列表 & 显示转码中、待审核等状态 \\
    \hline
  \end{tabular}
\end{table}

测试结果表明，创作中心能够完成视频投稿的基本流程。用户提交投稿后，不需要等待转码立即完成，可以通过稿件状态查看后续处理进度。这种测试结果符合前文对异步处理的设计要求。

\section{视频处理与审核流程测试}
\label{sec:video-audit-test}

\subsection{视频转码流程测试}
\label{subsec:transcode-test}

视频转码测试主要验证资源服务能否消费转码队列，并调用 FFmpeg 生成 HLS 播放资源。测试时，需要观察原始视频分片是否成功合并，转码后是否生成 \texttt{.m3u8} 文件和对应媒体分片，主站服务是否收到转码结果回写。

转码测试的关键点不只是文件是否生成，还包括视频状态是否正确变化。如果转码成功，视频应进入待审核状态；如果转码失败，系统应记录失败原因并更新状态。FFmpeg 在视频格式转换和媒体切片方面具有较强能力，因此适合作为转码测试的核心工具\cite{ffmpegdocs}。

表~\ref{tab:transcode-test} 展示了视频转码测试用例。

\begin{table}[htbp]
  \centering
  \caption{视频转码测试用例}
  \label{tab:transcode-test}
  \begin{tabular}{p{0.22\textwidth}p{0.34\textwidth}p{0.32\textwidth}}
    \hline
    \textbf{测试项} & \textbf{测试操作} & \textbf{预期结果} \\
    \hline
    队列消费 & 启动资源服务消费转码任务 & 能够从队列取出待转码文件 \\
    分片合并 & 合并上传分片 & 生成完整原始视频文件 \\
    视频转码 & 调用 FFmpeg 处理视频 & 生成 HLS 播放资源 \\
    状态回写 & 转码完成后通知主站服务 & 文件状态更新成功 \\
    异常处理 & 使用异常视频文件测试 & 状态更新为转码失败并记录原因 \\
    \hline
  \end{tabular}
\end{table}

测试结果表明，资源服务能够完成分片合并和转码处理，转码完成后也能够回写状态。对于异常视频文件，系统能够记录失败状态，避免视频继续进入审核和发布流程。

\subsection{AI 审核流程测试}
\label{subsec:ai-audit-test}

AI 审核流程测试主要验证审核任务能否正常创建、消息能否发送、AI 服务能否消费任务、审核结果能否回写，并在后台页面展示。由于 AI 审核流程涉及主站服务、RabbitMQ、AI 服务和数据库，测试时需要关注跨服务链路是否完整。

测试时，当视频转码完成并进入待审核状态后，系统应自动创建 AI 审核任务，并向审核请求队列发送消息。AI 服务消费任务后，执行视频内容分析，并把结果发送到结果队列。主站服务消费结果后，应更新审核任务状态和审核明细。RabbitMQ 的消息机制能够支持这类异步请求和结果回传流程\cite{rabbitmqdocs}。

表~\ref{tab:ai-audit-test} 展示了 AI 审核流程测试用例。

\begin{table}[htbp]
  \centering
  \caption{AI 审核流程测试用例}
  \label{tab:ai-audit-test}
  \begin{tabular}{p{0.22\textwidth}p{0.34\textwidth}p{0.32\textwidth}}
    \hline
    \textbf{测试项} & \textbf{测试操作} & \textbf{预期结果} \\
    \hline
    任务创建 & 视频转码完成后触发审核 & 生成审核任务和明细记录 \\
    消息发送 & 发送审核请求到 RabbitMQ & 请求队列中出现审核消息 \\
    AI 消费 & AI 服务消费审核请求 & 任务状态变为处理中 \\
    结果回写 & AI 服务返回审核结果 & 审核结果保存到数据库 \\
    后台展示 & 管理员查看审核详情 & 展示风险等级、标签和原因说明 \\
    \hline
  \end{tabular}
\end{table}

测试结果表明，AI 审核流程可以形成闭环。AI 服务能够对视频进行初步分析，后端能够接收并保存结构化结果，后台页面也能够展示审核摘要和风险信息。需要注意的是，AI 结果主要用于辅助人工审核，不能直接替代管理员判断。

\subsection{人工审核与发布测试}
\label{subsec:manual-audit-publish-test}

人工审核测试主要验证管理员能否查看待审核视频、进入审核详情、播放视频、查看 AI 审核结果，并执行通过或驳回操作。审核通过后，系统应将投稿数据迁移到正式视频表，并写入搜索索引；审核驳回后，系统应保留投稿状态和驳回原因。

人工审核测试需要特别关注正式发布数据和投稿数据之间的关系。如果审核通过后只更新状态而没有写入正式表，前台可能无法展示视频；如果未审核视频错误写入正式表，就会带来内容安全风险。

表~\ref{tab:manual-audit-test} 展示了人工审核与发布测试用例。

\begin{table}[htbp]
  \centering
  \caption{人工审核与发布测试用例}
  \label{tab:manual-audit-test}
  \begin{tabular}{p{0.22\textwidth}p{0.34\textwidth}p{0.32\textwidth}}
    \hline
    \textbf{测试项} & \textbf{测试操作} & \textbf{预期结果} \\
    \hline
    审核列表 & 管理员打开待审核列表 & 显示待审核视频和 AI 摘要 \\
    审核详情 & 查看视频详情和审核结果 & 视频、投稿信息和风险信息展示正常 \\
    审核通过 & 点击通过按钮 & 视频进入正式表并写入搜索索引 \\
    审核驳回 & 点击驳回并填写原因 & 投稿状态更新为审核不通过 \\
    前台验证 & 审核通过后搜索和播放视频 & 视频可被搜索并正常播放 \\
    \hline
  \end{tabular}
\end{table}

测试结果表明，后台人工审核能够完成最终发布控制。审核通过后，视频可以进入前台展示和播放；审核不通过时，视频不会进入正式数据集合。该结果说明视频内容发布流程具有基本安全性。

\section{异常场景测试}
\label{sec:exception-test}

\subsection{权限异常测试}
\label{subsec:permission-exception-test}

权限异常测试主要验证未登录用户、普通用户和管理员在访问不同接口时是否受到正确限制。未登录用户尝试评论、弹幕、点赞或投稿时，系统应提示登录；普通用户访问后台管理接口时，应被拒绝；外部请求访问内部接口时，网关应进行拦截。

权限控制测试不仅是安全测试，也是业务流程测试。因为很多功能依赖身份状态，如果权限控制不正确，可能导致游客产生互动数据，或者普通用户访问后台审核功能。Spring Security 等认证授权机制可以为权限设计提供支持\cite{springsecuritydocs}。

表~\ref{tab:permission-exception-test} 展示了权限异常测试用例。

\begin{table}[htbp]
  \centering
  \caption{权限异常测试用例}
  \label{tab:permission-exception-test}
  \begin{tabular}{p{0.24\textwidth}p{0.34\textwidth}p{0.30\textwidth}}
    \hline
    \textbf{测试场景} & \textbf{测试操作} & \textbf{预期结果} \\
    \hline
    未登录互动 & 游客发表评论或点赞 & 提示登录，不保存操作 \\
    普通用户访问后台 & 用户访问管理端接口 & 返回权限不足或拒绝访问 \\
    管理员未登录 & 直接访问后台审核页面 & 跳转登录或返回未授权 \\
    外部访问内部接口 & 请求包含内部接口前缀 & 网关拦截请求 \\
    \hline
  \end{tabular}
\end{table}

测试结果表明，系统能够对不同身份进行基本限制。未登录用户和普通用户无法执行越权操作，后台接口和内部接口能够得到保护。

\subsection{任务失败异常测试}
\label{subsec:task-failure-test}

任务失败异常测试主要验证视频转码失败、AI 审核失败和消息处理异常时系统是否能够记录状态。对于视频转码失败，系统应将稿件状态更新为转码失败；对于 AI 审核失败，系统应记录失败原因，并允许管理员后续人工处理。

异步任务最怕状态丢失。如果一个任务失败后没有任何记录，用户和管理员都不知道问题出在哪里。因此，任务状态和错误信息是异常测试中的重点。数据密集型应用设计中也强调，可靠系统需要能处理失败并保留恢复依据\cite{kleppmann2017ddia}。

表~\ref{tab:task-failure-test} 展示了任务失败异常测试用例。

\begin{table}[htbp]
  \centering
  \caption{任务失败异常测试用例}
  \label{tab:task-failure-test}
  \begin{tabular}{p{0.24\textwidth}p{0.34\textwidth}p{0.30\textwidth}}
    \hline
    \textbf{测试场景} & \textbf{测试操作} & \textbf{预期结果} \\
    \hline
    转码失败 & 上传异常格式视频或中断转码 & 状态更新为转码失败并记录原因 \\
    AI 请求失败 & 停止 RabbitMQ 或模拟消息发送异常 & 审核任务记录失败信息 \\
    AI 服务异常 & 停止 AI 服务后发送审核任务 & 请求保留或任务进入异常状态 \\
    结果回写失败 & 模拟结果消费异常 & 记录日志并保留消息处理依据 \\
    \hline
  \end{tabular}
\end{table}

测试结果表明，系统在异常情况下能够通过状态记录和错误信息降低流程中断影响。虽然毕业设计原型无法覆盖所有生产级异常场景，但已经具备基本的异常处理思路。

除了正常业务链路，系统还需要验证异常场景下的处理效果。异常测试的重点不是让系统永远不出错，而是在错误发生时能够返回合理提示、记录失败状态，并避免业务数据混乱。算法~\ref{alg:test-exception-regression} 展示了异常场景回归测试流程。

\begin{algorithm}[htbp]
  \caption{异常场景回归测试流程}
  \label{alg:test-exception-regression}
  \begin{algorithmic}[1]
    \Require 异常测试用例集合 $caseList$
    \Ensure 异常测试结果
    \For{每个测试用例 $case$ in $caseList$}
      \State 初始化测试数据和系统状态
      \If{$case$ 为未登录互动}
        \State 游客调用评论、弹幕或点赞接口
        \State 校验系统是否返回未登录提示且未写入行为数据
      \ElsIf{$case$ 为转码失败}
        \State 上传异常视频文件或模拟 FFmpeg 失败
        \State 校验投稿状态是否更新为“转码失败”
      \ElsIf{$case$ 为 AI 服务不可用}
        \State 停止 AI 服务后发送审核任务
        \State 校验审核任务是否保留异常状态或等待处理
      \ElsIf{$case$ 为普通用户访问后台}
        \State 使用普通用户令牌访问后台审核接口
        \State 校验系统是否拒绝访问
      \ElsIf{$case$ 为外部访问内部接口}
        \State 直接请求包含 $/innerApi$ 的路径
        \State 校验网关是否拦截请求
      \EndIf
      \State 记录测试结果和异常日志
    \EndFor
  \end{algorithmic}
\end{algorithm}

由算法~\ref{alg:test-exception-regression} 可以看出，异常测试覆盖了权限异常、转码异常、AI 服务异常和内部接口访问异常等情况。这些测试有助于验证系统在非正常情况下是否仍能保持状态清晰，并为后续问题定位提供依据。


\section{测试结果分析}
\label{sec:test-result-analysis}

\subsection{功能测试结果分析}
\label{subsec:function-test-result}

从功能测试结果看，系统能够完成普通用户端的主要操作，包括登录、浏览、搜索、播放、评论、弹幕、点赞、收藏和投币等功能。创作中心能够完成封面上传、视频分片上传、投稿提交和稿件状态查看。后台管理端能够完成管理员登录、分类管理、审核列表、审核详情和人工审核操作。

视频处理流程测试表明，资源服务可以从队列中获取待转码文件，完成分片合并和 HLS 资源生成。审核流程测试表明，主站服务能够创建 AI 审核任务，AI 服务能够返回审核结果，后台页面能够展示相关风险信息。

这些结果说明，系统已经实现在线视频分享平台的核心业务闭环。简单来说，用户能上传视频，系统能处理视频，管理员能审核视频，审核通过后普通用户能搜索和播放视频。

\subsection{非功能测试结果分析}
\label{subsec:nonfunction-test-result}

从非功能角度看，系统在安全性、可靠性和可维护性方面具备基本设计。安全性方面，系统能够区分游客、普通用户和管理员；可靠性方面，转码失败和 AI 审核失败都有状态记录；可维护性方面，系统通过服务拆分和模块划分降低了功能混杂问题。

不过，系统仍属于毕业设计原型，在大规模并发、自动化测试覆盖率、分布式文件存储和模型审核准确率评估方面还有提升空间。比如当前文件存储更适合本地演示，如果后续要部署到多实例环境，就需要迁移到对象存储或统一文件服务。

表~\ref{tab:test-result-summary} 对系统测试结果进行了汇总。

\begin{table}[htbp]
  \centering
  \caption{系统测试结果汇总}
  \label{tab:test-result-summary}
  \begin{tabular}{p{0.22\textwidth}p{0.34\textwidth}p{0.34\textwidth}}
    \hline
    \textbf{测试类别} & \textbf{测试内容} & \textbf{结果分析} \\
    \hline
    用户端功能 & 登录、浏览、搜索、播放、互动 & 能够覆盖普通用户主要使用场景 \\
    创作中心 & 分片上传、投稿提交、稿件管理 & 能够完成视频投稿和状态查看流程 \\
    资源处理 & 分片保存、转码、HLS 输出 & 支持视频从原始文件到网页播放资源的转换 \\
    后台管理 & 分类管理、视频审核、审核详情 & 支持管理员进行内容管理和人工审核 \\
    AI 审核 & 审核任务、消息通信、结果展示 & 支持 AI 初审和人工终审的协同流程 \\
    异常场景 & 权限不足、转码失败、AI 失败 & 能够通过状态记录和人工处理避免流程中断 \\
    \hline
  \end{tabular}
\end{table}

由表~\ref{tab:test-result-summary} 可知，系统主要功能测试结果符合预期。系统能够完成从用户投稿到视频发布的核心流程，同时能够对异常情况进行基本处理。

\section{本章小结}
\label{sec:chapter6-summary}

本章对系统测试过程和测试结果进行了说明。首先，介绍了系统测试环境和测试方案，明确了测试需要覆盖用户端、创作中心、后台管理、视频处理、AI 审核和异常场景。随后，围绕用户认证、首页搜索、视频播放、互动功能、投稿上传、视频转码、AI 审核和人工审核等功能设计了测试用例。

测试结果表明，系统能够完成在线视频分享平台的主要业务流程。用户可以浏览、搜索和播放视频，创作者可以上传和投稿视频，资源服务可以完成视频转码，后台管理员可以进行审核，AI 服务可以提供初步风险分析。异常测试也表明，系统能够对未登录访问、权限不足、转码失败和 AI 审核失败等情况进行基本处理。

总体来看，本系统已经满足毕业设计对在线视频分享平台原型的核心要求。不过，系统在高并发性能测试、自动化测试覆盖率、文件存储扩展和 AI 审核准确率评估等方面仍有继续优化空间。

\chapter{总结与展望}
\label{chap:conclusion-outlook}

\section{研究工作总结}
\label{sec:research-summary}

本文围绕“基于微服务架构的在线视频分享平台设计与实现”展开研究与开发，完成了一个具有视频浏览、视频投稿、分片上传、转码播放、用户互动、后台管理和 AI 审核能力的在线视频分享平台原型。系统采用前后端分离架构，后端按照业务能力拆分为网关服务、主站服务、资源服务、互动服务、后台服务和公共基础模块，AI 审核服务则以独立服务形式运行。

在需求分析阶段，本文从游客、注册用户、内容创作者、管理员和 AI 审核服务等角色出发，分析了系统需要实现的功能需求和非功能需求。普通用户需要完成视频浏览、搜索、播放和互动；内容创作者需要完成视频投稿和稿件管理；管理员需要完成分类管理和视频审核；AI 审核服务需要对上传视频进行自动化风险分析。系统还需要满足可维护性、可扩展性、安全性和数据一致性等非功能需求。

在系统设计阶段，本文对系统总体架构、微服务模块划分、网关与接口、数据库、缓存搜索、文件存储、消息通信和核心业务流程进行了设计。系统通过网关统一入口，通过 Nacos 管理服务发现，通过 OpenFeign 完成部分同步调用，通过 Redis 保存缓存和转码队列，通过 RabbitMQ 完成 AI 审核消息通信，通过 MySQL 保存核心数据，通过 Elasticsearch 支撑视频搜索。这样的设计体现了微服务架构在复杂业务拆分中的应用价值\cite{lewis2014microservices}。

在系统实现阶段，本文重点完成了视频生命周期相关功能。用户在创作中心上传视频和封面后，系统采用分片上传方式保存视频文件，并在投稿提交后将待转码任务写入队列。资源服务后台处理转码任务，调用 FFmpeg 对视频进行合并、转码和 HLS 切片。主站服务在确认所有分 P 转码完成后，将视频状态更新为待审核，并创建 AI 审核任务。AI 服务消费审核请求后，对视频进行音频提取、语音转写、声音事件识别、关键帧抽取和多模态分析，并将审核结果返回后端。管理员最终在后台结合 AI 结果执行人工通过或驳回。

在系统测试阶段，本文围绕用户端、创作中心、资源处理、后台管理、AI 审核和异常场景设计了测试用例。测试内容覆盖登录、视频列表、视频播放、评论弹幕、点赞收藏投币、封面上传、视频分片上传、视频转码、审核任务创建、AI 结果回写和后台人工审核等功能。测试结果表明，系统能够完成在线视频分享平台的主要业务流程，并能够通过异步任务和服务拆分提高系统结构清晰度。

\section{项目创新点}
\label{sec:project-innovation}

本项目的创新点主要体现在三个方面。第一，系统不是只实现简单的视频播放，而是围绕视频生命周期构建了从投稿、上传、转码、审核、发布到播放互动的完整流程。相比普通信息管理系统，该项目更接近真实视频平台的业务链路。

第二，系统将微服务架构与视频资源处理结合起来。主站服务、资源服务、互动服务、后台服务和 AI 服务之间职责相对清晰。资源服务专门处理视频分片上传和转码，互动服务专门处理评论、弹幕和用户行为，AI 服务专门处理内容审核。这种拆分方式降低了不同业务之间的耦合度。

第三，系统引入了 AI 辅助审核流程。AI 审核服务通过消息队列消费审核任务，对视频语音、声音事件和画面内容进行综合分析，并返回风险等级、风险标签和命中片段。管理员再结合 AI 审核结果进行人工终审。该设计体现了人机协同审核思想，也提高了系统内容治理能力。多模态理解模型的发展为视频内容审核提供了更丰富的技术基础\cite{bai2023qwen}。

表~\ref{tab:project-innovation} 对项目主要创新点进行了总结。

\begin{table}[htbp]
  \centering
  \caption{项目主要创新点}
  \label{tab:project-innovation}
  \begin{tabular}{p{0.24\textwidth}p{0.66\textwidth}}
    \hline
    \textbf{创新点} & \textbf{具体说明} \\
    \hline
    完整视频生命周期 & 实现视频从投稿、上传、转码、审核、发布到播放互动的闭环流程 \\
    微服务业务拆分 & 将主站、资源、互动、后台和 AI 审核能力拆分为相对独立模块 \\
    AI 辅助审核 & 结合语音识别、声音事件识别和多模态理解辅助管理员审核 \\
    人机协同机制 & AI 负责初审提示，管理员负责最终审核，兼顾效率和准确性 \\
    \hline
  \end{tabular}
\end{table}

由表~\ref{tab:project-innovation} 可以看出，本项目的重点不是某一个单独功能，而是多个功能之间的协同。视频处理、微服务拆分和 AI 审核共同构成了项目的核心特点。

\section{项目不足}
\label{sec:project-limitations}

虽然本系统已经完成了在线视频分享平台的核心功能，但由于时间、硬件环境和项目规模限制，仍然存在一些不足。第一，系统目前更偏向毕业设计原型，尚未进行大规模并发测试。在真实生产环境中，视频平台会面对大量用户访问、视频上传和播放请求，系统还需要进一步进行性能压测和容量规划。

第二，文件存储方式仍以本地目录为主。该方式适合本地开发和演示，但如果部署到多实例服务环境中，就会出现文件共享和访问路径管理问题。后续可以考虑接入对象存储或分布式文件系统，提高文件存储的扩展性。

第三，AI 审核准确率还需要进一步验证。当前 AI 审核流程能够输出风险提示和审核建议，但模型结果可能受到视频内容、模型能力和提示词设计影响。对于复杂语义、隐晦表达或低质量视频，AI 仍可能出现误判。因此，当前 AI 审核更适合作为辅助工具，而不是完全自动化审核。

第四，系统自动化测试覆盖还不够完整。当前测试主要围绕核心功能和主要流程展开，尚未形成完整的单元测试、接口自动化测试和持续集成流程。软件质量模型强调可靠性和可维护性，自动化测试是后续提升系统质量的重要方向\cite{iso25010}。

表~\ref{tab:project-limitations} 对系统不足进行了整理。

\begin{table}[htbp]
  \centering
  \caption{系统不足分析}
  \label{tab:project-limitations}
  \begin{tabular}{p{0.24\textwidth}p{0.36\textwidth}p{0.30\textwidth}}
    \hline
    \textbf{不足类别} & \textbf{具体表现} & \textbf{影响} \\
    \hline
    并发测试不足 & 未进行大规模用户访问压测 & 难以评估生产环境承载能力 \\
    文件存储限制 & 主要使用本地文件目录 & 多实例部署时文件共享不方便 \\
    AI 准确率不足 & 模型误判和漏判仍可能存在 & 审核结果仍需人工确认 \\
    自动化测试不足 & 缺少完整接口测试和持续集成 & 后续迭代回归成本较高 \\
    推荐能力不足 & 当前更偏向搜索和基础展示 & 个性化内容分发能力有限 \\
    \hline
  \end{tabular}
\end{table}

由表~\ref{tab:project-limitations} 可以看出，系统当前已经具备核心功能，但距离商业级视频平台仍有差距。这些不足也为后续优化提供了方向。

\section{后续优化方向}
\label{sec:future-work}

针对当前系统存在的不足，后续可以从以下几个方面进行优化。第一，优化视频资源存储方式。可以将本地文件存储迁移到对象存储或分布式文件系统，并结合 CDN 提高视频访问速度。这样可以解决多实例部署时文件共享不便的问题，也能提高视频加载效率。

第二，提升系统性能和并发能力。后续可以对首页列表、视频播放、评论弹幕、搜索接口和上传接口进行压力测试，并根据测试结果优化缓存策略、数据库索引和服务扩容方式。对于播放量、点赞数等高频统计数据，可以进一步使用异步队列和批量更新减少数据库压力。

第三，完善 AI 审核能力。后续可以引入更完善的审核规则体系，结合模型结果、关键词规则和人工反馈持续优化审核效果。也可以对不同类型风险建立更细的标签体系，例如暴力、低俗、敏感语音、异常声音等，让审核结果更容易理解。算法内容审核研究也指出，自动化审核需要与人工治理和规则体系结合，才能更好地服务平台管理\cite{gorwa2020algorithmic}。

第四，增强平台推荐和数据分析能力。当前系统主要实现视频浏览和搜索，后续可以基于用户观看、点赞、收藏、投币和关注等行为数据，构建简单推荐策略。比如根据用户关注、分类偏好和热门视频进行推荐，让用户更容易发现感兴趣的内容。

第五，完善测试和部署体系。后续可以补充单元测试、接口自动化测试和端到端测试，并结合持续集成工具进行自动构建和部署。对于微服务系统，还可以增加服务监控、日志追踪和异常告警，提高系统运行可观察性。

表~\ref{tab:future-work} 总结了系统后续优化方向。

\begin{table}[htbp]
  \centering
  \caption{后续优化方向}
  \label{tab:future-work}
  \begin{tabular}{p{0.24\textwidth}p{0.42\textwidth}p{0.24\textwidth}}
    \hline
    \textbf{优化方向} & \textbf{优化内容} & \textbf{预期效果} \\
    \hline
    文件存储优化 & 接入对象存储、分布式文件系统或 CDN & 提高视频访问和部署扩展能力 \\
    性能优化 & 压力测试、缓存优化、索引优化、服务扩容 & 提升高并发访问能力 \\
    AI 审核优化 & 完善风险标签、规则体系和人工反馈机制 & 提高审核准确率和可解释性 \\
    推荐能力增强 & 基于用户行为和分类偏好进行推荐 & 提升内容分发效果 \\
    测试部署完善 & 增加自动化测试、日志追踪和监控告警 & 降低后续维护成本 \\
    \hline
  \end{tabular}
\end{table}

由表~\ref{tab:future-work} 可以看出，后续优化不仅包括功能扩展，也包括性能、部署、审核和测试体系的完善。对于一个在线视频分享平台来说，功能完整只是第一步，后续更重要的是提高系统稳定性、扩展性和内容治理能力。

\section{本章小结}
\label{sec:chapter7-summary}

本章对全文工作进行了总结，并分析了项目创新点、不足和后续优化方向。总体来看，本文完成了一个基于微服务架构的在线视频分享平台原型，实现了视频浏览、视频投稿、分片上传、转码播放、互动行为、后台审核和 AI 辅助审核等功能。系统通过服务拆分、异步任务和人机协同审核机制，形成了较完整的视频内容处理闭环。

同时，系统仍存在一些不足，例如并发性能尚未充分验证、文件存储方式仍偏本地化、AI 审核准确率有待进一步评估、自动化测试体系还不够完善等。后续可以从对象存储、性能优化、审核规则完善、推荐算法和自动化部署等方面继续改进。

通过本课题的设计与实现，可以较完整地理解在线视频分享平台的业务流程和技术结构，也能够体现微服务架构、视频处理、异步任务和智能审核在复杂 Web 系统中的综合应用价值。
